\begin{question}
\label{ex:mod3M}
Consider the following scenario.  Each thread has an integer-valued identity.
A particular resource should be accessed according to the following
constraint: a new thread can start using the resource only if the sum of the
identities of those threads currently using it is divisible by~3.  Implement a
JVM monitor that enforces this, extending the following trait.
\begin{scala}
trait Mod3{
  def enter(id: Int): Unit
  def exit(id: Int): Unit
}
\end{scala}

Test your program using the linearizability framework.  \textbf{Hint:} one
approach is to use an immutable |Set[Int]| as the sequential specification
object, representing the set of identities of threads currently using the
resource, and to include a suitable precondition for the sequential version of
|enter|.
\end{question}

%%%%%

\begin{answerI}
Here's my code.
%
\begin{scala}
class Mod3M{
  /** Sum of identites of threads currently in critical region. */
  private var current = 0

  /** Enter the critical region. */
  def enter(id: Int) = synchronized{
    while(current%3 != 0) wait() 
    current += id
    if(current%3 == 0) notify()
  }

  /** Leave the critical region. */
  def exit(id: Int) = synchronized{
    current -= id
    if(current%3 == 0) notify()
  }
}
\end{scala}
%
Note that an exiting thread performs a |notify| only if there is some point in
doing so, i.e., |current| is divisible by~3.  It would be ok to do a
|notifyAll| at this point; however, that is quite likely to wake up several
threads, and often only one of them will be able to enter.  Instead, it's more
efficient to do a |notify|.  However, this means that the entering thread
should maybe signal to another waiting thread. 


%% Alternatively, one can note that the following is invariant: at most one
%% thread whose identity is not divisible by 3 is present at any time.
%% Therefore it would be enough to store a boolean recording whether such a
%% thread is currently present.

For testing, we can follow the hint and use a |Set[Int]| as the specification
object.  The sequential operations are then as follows.
%
\begin{scala}
  type S = Set[Int]
  def seqEnter(id: Int)(current: S): (Unit, S) = {
    require(current.sum%3 == 0); ((), current+id)
  }
  def seqExit(id: Int)(current: S): (Unit, S) = {
    assert(current.contains(id)); ((), current-id)
  }
\end{scala}
%
In particular, the |seqEnter| operation captures the precondition.  (The
|assert| in |seqExit| is simply a sanity check.)  We can then run threads that
repeatedly enter and exit.  The rest of the tester is standard. 

Alternatively, the testing could be based on a sequential specification object
that just records the \emph{sum} of the identities of current threads.
%
\begin{scala}
  type S = Int
  def seqEnter(id: Int)(current: S): (Unit, S) = {
    require(current%3 == 0); ((), current+id)
  }

  def seqExit(id: Int)(current: S): (Unit, S) = {
    assert(id <= current); ((), current-id)
  }
\end{scala}
\end{answerI}
